% ============================================================
%  EXAMEN TEMA 4: ECUACIONES DE PRIMER GRADO EN ℤ
%  Curso Introductorio de Matemáticas — 1er Año
%  Autor: Ing. Gabriel Astudillo
%  Sesiones cubiertas: Semana 4 (clases 13 a 16)
%  Nota: enunciado limpio (sin pistas ni desarrollo). Las
%  soluciones están en la "Clave de Respuestas" al final.
% ============================================================

\documentclass[12pt, letterpaper]{article}

% ── Paquetes ─────────────────────────────────────────────────

\usepackage{mathwizards-guia}
\mwguia{Examen — Tema 4: Ecuaciones de Primer Grado}{Ing. Gabriel Astudillo $\cdot$ Curso 1er Año}

% ── Comandos propios de este examen ───────────────────
\newcommand{\Z}{\mathbb{Z}}

\begin{document}
% ============================================================

% ── Portada / Título ─────────────────────────────────────────
\mwportada{Examen del Tema 4}{Ecuaciones de Primer Grado en $\Z$}{Igualdad, Lenguaje Algebraico y Resolución Paso a Paso}

\vspace{4pt}
\begin{cuadroinfo}
  \small
  \textbf{Nombre:} \rule{0.55\linewidth}{0.4pt} \quad \textbf{Fecha:} \rule{0.15\linewidth}{0.4pt} \\[6pt]
  \textbf{Tiempo sugerido:} 60 minutos \quad$\bullet$\quad \textbf{Puntaje total:} 100 puntos \\[4pt]
  \textbf{Instrucciones:}
  \begin{enumerate}[leftmargin=*]
    \item Lee cada ejercicio con calma antes de resolverlo.
    \item Muestra \emph{todos} tus pasos: se evalúa cómo llegas a la respuesta, no solo la respuesta.
    \item En las ecuaciones, \textbf{comprueba} siempre tu solución sustituyendo en la igualdad original.
    \item Con paréntesis o corchetes, \textbf{elimina primero los signos de agrupación} y luego resuelve.
    \item No se permite el uso de calculadora.
  \end{enumerate}
\end{cuadroinfo}

\vspace{8pt}

% ============================================================
\section{Lenguaje Algebraico y Planteamiento}
% ============================================================

\begin{enumerate}[leftmargin=*]
  \item Escriba en lenguaje algebraico:
        \begin{enumerate}[label=\alph*)]
          \item El triple de un número, disminuido en $15$.
          \item El doble de la diferencia entre un número y $9$.
          \item La mitad de la suma entre un número y $7$, aumentada en $5$.
        \end{enumerate}

  \item Identifique el primer miembro, el segundo miembro, la incógnita, los términos y el coeficiente de la ecuación:
        \[
          -5x + 12 = -18
        \]

  \item Un número, aumentado en $4$, es igual al doble de su diferencia con $3$. Plantee la ecuación (no la resuelva).

  \item El perímetro de un rectángulo es $54$ cm y su largo mide $3$ cm más que el doble de su ancho. Plantee la ecuación (no la resuelva).
\end{enumerate}

\newpage

% ============================================================
\section{Pasos para Resolver y Comprobar}
% ============================================================

\begin{enumerate}[leftmargin=*]
  \item Resuelva y compruebe:
        \[
          6x - 15 = 2x + 25
        \]

  \item Resuelva y compruebe:
        \[
          -3x + 14 = -2x - 9
        \]

  \item Resuelva aplicando operaciones inversas:
        \[
          8 - 5x = 2x - 13
        \]

  \item Resuelva juntando primero los términos semejantes de cada miembro:
        \[
          5x - 9 + 2x = 3x + 15
        \]
\end{enumerate}

\newpage

% ============================================================
\section{Ecuaciones Sin Signos de Agrupación}
% ============================================================

\begin{enumerate}[leftmargin=*]
  \item Resuelva:
        \[
          9x - 14 = 4x + 21
        \]

  \item Resuelva:
        \[
          4x - 18 + 6x = 15 - 3x + 19
        \]

  \item Resuelva y determine si la solución pertenece a $\Z$:
        \[
          7x - 18 = 3x + 22
        \]

  \item Resuelva y compruebe:
        \[
          -5x + 18 - 2x = 4 - 9x + 26
        \]
\end{enumerate}

\newpage

% ============================================================
\section{Ecuaciones con Paréntesis y Problemas de la Vida Real}
% ============================================================

\begin{enumerate}[leftmargin=*]
  \item Resuelva:
        \[
          2[4 - 3(x - 2)] = 4(x + 5)
        \]

  \item Resuelva:
        \[
          5(2 - x) - 3[2x - (x + 1)] = -19
        \]

  \item \textbf{Problema de edades:} Un padre tiene $42$ años y su hija $12$ años. \textquestiondown Dentro de cuántos años la edad del padre será el doble de la de su hija?

  \item \textbf{Problema de adivinar el número:} «Pensé un número entero, lo multipliqué por $-3$, al resultado le sumé $18$, luego multipliqué todo por $2$ y finalmente le resté $10$, obteniendo $32$». \textquestiondown Cuál fue el número que pensé?
\end{enumerate}

\newpage

% ============================================================
\ifmwclaves
  \section{Clave de Respuestas}
  % ============================================================

  \subsection*{Lenguaje Algebraico y Planteamiento (Ejercicios 1--4)}

  \begin{enumerate}[leftmargin=*, label=\textbf{\arabic*.}]
    \item a) $3x - 15$ \quad b) $2(x - 9)$ \quad c) $\dfrac{x + 7}{2} + 5$
    \item 1er miembro: $-5x + 12$; 2do miembro: $-18$; incógnita: $x$; términos: $-5x$, $12$, $-18$; coeficiente de $x$: $-5$
    \item $x + 4 = 2(x - 3)$
    \item Sea $a$ el ancho; el largo es $2a + 3$. $2(a + 2a + 3) = 54$
  \end{enumerate}

  \subsection*{Pasos para Resolver y Comprobar (Ejercicios 1--4)}

  \begin{enumerate}[leftmargin=*, label=\textbf{\arabic*.}]
    \item $6x - 2x = 25 + 15 \Rightarrow 4x = 40 \Rightarrow x = 10$ \ ($6(10) - 15 = 45$; $2(10) + 25 = 45$ \checkmark)
    \item $-3x + 2x = -9 - 14 \Rightarrow -x = -23 \Rightarrow x = 23$
    \item $8 + 13 = 2x + 5x \Rightarrow 21 = 7x \Rightarrow x = 3$ \ ($8 - 15 = -7$; $6 - 13 = -7$ \checkmark)
    \item $7x - 9 = 3x + 15 \Rightarrow 4x = 24 \Rightarrow x = 6$
  \end{enumerate}

  \subsection*{Ecuaciones Sin Signos de Agrupación (Ejercicios 1--4)}

  \begin{enumerate}[leftmargin=*, label=\textbf{\arabic*.}]
    \item $9x - 4x = 21 + 14 \Rightarrow 5x = 35 \Rightarrow x = 7$
    \item $10x - 18 = 34 - 3x \Rightarrow 13x = 52 \Rightarrow x = 4$
    \item $7x - 3x = 22 + 18 \Rightarrow 4x = 40 \Rightarrow x = 10$ \ $\Rightarrow$ sí pertenece a $\Z$
    \item $-5x + 18 - 2x = 4 - 9x + 26 \Rightarrow -7x + 18 = 30 - 9x \Rightarrow 2x = 12 \Rightarrow x = 6$ \ ($-5(6)+18-2(6) = -24$; $4-9(6)+26 = -24$ \checkmark)
  \end{enumerate}

  \subsection*{Ecuaciones con Paréntesis y Problemas de la Vida Real (Ejercicios 1--4)}

  \begin{enumerate}[leftmargin=*, label=\textbf{\arabic*.}]
    \item $4 - 3(x - 2) = -3x + 10$; $2(-3x + 10) = 4x + 20 \Rightarrow -6x + 20 = 4x + 20 \Rightarrow x = 0$
    \item $2x - (x + 1) = x - 1$; $\Rightarrow 5(2 - x) - 3(x - 1) = -19 \Rightarrow 10 - 5x - 3x + 3 = -19 \Rightarrow 13 - 8x = -19 \Rightarrow -8x = -32 \Rightarrow x = 4$
    \item Sea $t$ los años. $42 + t = 2(12 + t) \Rightarrow 42 + t = 24 + 2t \Rightarrow t = 18$ años
    \item Ecuación: $2(-3x + 18) - 10 = 32 \Rightarrow -6x + 36 - 10 = 32 \Rightarrow -6x + 26 = 32 \Rightarrow -6x = 6 \Rightarrow x = -1$
  \end{enumerate}

  \vspace{10pt}
  \begin{cuadrosolucion}
    \textbf{\textquestiondown Cómo te fue?} Si comprobaste todas tus soluciones, vas por muy buen camino. Recuerda: el \emph{lenguaje algebraico} es la base de todo, y \emph{comprobar} es el sello de un buen resolutor. \textquestiondown Listo para el examen final!
  \end{cuadrosolucion}

\fi
\end{document}
